\UPSTIactiviteStyleCadre
\section{Fiche de synthèse}

Il est temps de réaliser une fiche de synthèse résumant tous les points importants vus depuis le début de l'année.

\begin{UPSTIactivite}{Fiche de synthèse (30 min en classe, à terminer à la maison.)}
    \UPSTIetape{\textbf{Au brouillon} : Lister les titres des notions vues depuis le début de l'année.}
    \UPSTIetape{\textbf{Au propre} : Pour chaque notion, rédiger une phrase résumant l'essentiel à retenir.}

    Voici les éléments indispensables pour chaque fiche de synthèse :
    \begin{itemize}
        \item Titre de la fiche de synthèse (ex : \og Les boucles en C \fg)
        \item Définition de la notion
        \item Utilité de la notion
        \item Syntaxe en langage C
        \item Exemple d'utilisation
    \end{itemize}

    \textbf{Faire vérifier la fiche de synthèse par l'enseignant.}
\end{UPSTIactivite}
